\chapter{Tratamento de Erros}

\section{Modelo de erros}

\hayashi{} distingue dois tipos de erro:

\begin{description}
  \item[Erros de parse] Detectados antes da execução, ao analisar o código.
        Incluem linha e trecho do código.
  \item[Erros de runtime] Ocorrem durante a execução — variável indefinida,
        tipo errado, divisão por zero, chave ausente.
\end{description}

Ambos os tipos exibem o trecho de código e a linha, facilitando a localização:

\begin{lstlisting}[style=inline]
error: line 5: undefined variable 'rendimento' — did you mean 'rendimento_bruto'?
   5 │ display rendimento * 12
     │ ^^^^^^^^^^^^^^^^^^^^^^^^
\end{lstlisting}

\section{\texttt{try} / \texttt{catch}}

Erros de runtime podem ser capturados com \lstinline{try}/\lstinline{catch}:

\begin{lstlisting}[caption={try/catch básico}]
try {
  let resultado = 10 / 0
  display resultado
} catch e {
  display f"Erro capturado: {e}"
}
\end{lstlisting}

A variável \lstinline{e} dentro do \lstinline{catch} contém a mensagem de erro
como string.

\subsection*{Capturando erros de arquivo}

\begin{lstlisting}
try {
  load "dados_inexistentes.csv" as df
} catch e {
  display f"Arquivo não encontrado: {e}"
}
\end{lstlisting}

\subsection*{\texttt{try} como expressão}

\lstinline{try}/\lstinline{catch} pode retornar um valor:

\begin{lstlisting}
let valor = try {
  int("não é um número")
} catch _ {
  -1
}
display valor    // -1
\end{lstlisting}

\section{Erros comuns e suas mensagens}

\begin{center}
\begin{tabular}{ll}
\toprule
\textbf{Situação} & \textbf{Mensagem típica} \\
\midrule
Variável não declarada   & \texttt{undefined variable 'x' — did you mean 'y'?} \\
Função não existe        & \texttt{undefined function 'foo' — did you mean 'for'?} \\
Tipo errado              & \texttt{type mismatch — expected Int, got String} \\
Índice fora de limite    & \texttt{index 5 out of bounds (len=3)} \\
Chave ausente em dict    & \texttt{key 'nome' not found in dict} \\
Divisão por zero (Int)   & \texttt{division by zero} \\
Arquivo não encontrado   & \texttt{file not found: 'caminho.csv'} \\
\bottomrule
\end{tabular}
\end{center}

\section{Sugestões automáticas}

Quando o nome de uma variável ou função é parecido com algo existente, \hayashi{}
sugere a correção:

\begin{lstlisting}[style=inline]
error: undefined function 'regres' — did you mean 'regress'?
\end{lstlisting}

As sugestões usam distância de edição (\textit{Levenshtein}) sobre uma lista de
todos os nomes definidos no escopo e nas funções embutidas.

\section{Re-lançamento de erros}

Para propagar um erro capturado:

\begin{lstlisting}
try {
  operacao_arriscada()
} catch e {
  display f"log: {e}"
  error(e)    // re-lança o erro
}
\end{lstlisting}

A função \lstinline{error(msg)} lança um erro de runtime com a mensagem fornecida.
